%%%%%%%%%%%%%%%%%%%%%%%%%%%%%%%%%%%%%%%%%%%%%%%%%%%%%%%%%%%%%%%%%%%%%%%%%%%%%%%
% CHAPTER 6: THE REFERENCE STRUCTURE C*
%%%%%%%%%%%%%%%%%%%%%%%%%%%%%%%%%%%%%%%%%%%%%%%%%%%%%%%%%%%%%%%%%%%%%%%%%%%%%%%
% Version: 2.0 (January 2026)
%
% Purpose: Conceptual introduction to C* as the self-consistent reference
%          structure. All formal content (axioms, theorems, proofs) is in
%          Appendix UNMAPPED_PRF (FORMAL-CSTAR).
%
% Primary Appendix: D - FORMAL-CSTAR (Complete C* Theory)
% Secondary Appendices: B (Complementarity), BBB (Calibration)
%
% Prerequisites: Ch. 5 (Complementarity)
% Leads to: Ch. 7 (Coherence Index), Ch. 9 (Context)
%
% Chapter Type: B (Foundation Chapter - conceptual overview)
%%%%%%%%%%%%%%%%%%%%%%%%%%%%%%%%%%%%%%%%%%%%%%%%%%%%%%%%%%%%%%%%%%%%%%%%%%%%%%%

\chapter{The Reference Structure $C^*$}
\label{ch:reference-structure}

%%%%%%%%%%%%%%%%%%%%%%%%%%%%%%%%%%%%%%%%%%%%%%%%%%%%%%%%%%%%%%%%%%%%%%%%%%%%%%%
% CHAPTER SCOPE BOX
%%%%%%%%%%%%%%%%%%%%%%%%%%%%%%%%%%%%%%%%%%%%%%%%%%%%%%%%%%%%%%%%%%%%%%%%%%%%%%%
\begin{tcolorbox}[
    colback=purple!5!white,
    colframe=purple!75!black,
    title={\textbf{Chapter Scope: Ziel / In-Scope / Out-of-Scope}},
    fonttitle=\bfseries
]

\textbf{Ziel (Objective):}
Develop the conceptual understanding of reference points as a foundation for the EBF framework.

\vspace{0.3cm}
\textbf{In-Scope:}
\begin{itemize}[nosep]
    \item Conceptual introduction to reference points
    \item Motivation and intuition
    \item Connection to subsequent chapters
    \item Key definitions and concepts
\end{itemize}

\vspace{0.3cm}
\textbf{Out-of-Scope (delegiert an Appendices):}
\begin{itemize}[nosep]
    \item Complete formal treatment $\rightarrow$ FORMAL appendices
    \item Empirical validation $\rightarrow$ METHOD appendices
    \item Domain applications $\rightarrow$ Application chapters
\end{itemize}

\end{tcolorbox}



%%%%%%%%%%%%%%%%%%%%%%%%%%%%%%%%%%%%%%%%%%%%%%%%%%%%%%%%%%%%%%%%%%%%%%%%%%%%%%%
% QUICK REFERENCE BOX
%%%%%%%%%%%%%%%%%%%%%%%%%%%%%%%%%%%%%%%%%%%%%%%%%%%%%%%%%%%%%%%%%%%%%%%%%%%%%%%
\begin{tcolorbox}[
    colback=gray!5!white,
    colframe=gray!75!black,
    title={\textbf{Begriffe in diesem Kapitel}},
    fonttitle=\bfseries
]
\small
\begin{tabular}{@{}ll@{}}
\textbf{Key Concept} & \textbf{Definition} \\
\midrule
Reference points & See Section \ref{sec:ch6-intro} \\
\end{tabular}
\end{tcolorbox}



%%%%%%%%%%%%%%%%%%%%%%%%%%%%%%%%%%%%%%%%%%%%%%%%%%%%%%%%%%%%%%%%%%%%%%%%%%%%%%%
% INTUITION BOX
%%%%%%%%%%%%%%%%%%%%%%%%%%%%%%%%%%%%%%%%%%%%%%%%%%%%%%%%%%%%%%%%%%%%%%%%%%%%%%%
\begin{tcolorbox}[
    colback=blue!5!white,
    colframe=blue!75!black,
    title={\textbf{Intuition: A Motivating Example}},
    fonttitle=\bfseries
]
Consider \textbf{Anna}, a professional facing a decision that involves reference points.

She initially evaluates the choice using standard criteria, but something feels incomplete.
\textbf{Thomas}, her colleague, suggests considering additional dimensions beyond the obvious.

This chapter explores what Anna discovers when she applies the EBF framework---how
reference points interact with broader welfare considerations.

\medskip
\textit{By the end of this chapter, you will understand why Anna's initial analysis
was incomplete and how the EBF approach provides a more comprehensive view.}
\end{tcolorbox}



%%%%%%%%%%%%%%%%%%%%%%%%%%%%%%%%%%%%%%%%%%%%%%%%%%%%%%%%%%%%%%%%%%%%%%%%%%%%%%%
% CENTRAL QUESTION BOX
%%%%%%%%%%%%%%%%%%%%%%%%%%%%%%%%%%%%%%%%%%%%%%%%%%%%%%%%%%%%%%%%%%%%%%%%%%%%%%%
\begin{tcolorbox}[
    colback=red!5!white,
    colframe=red!75!black,
    title={\textbf{Central Question}},
    fonttitle=\bfseries
]
\Large
\centering
\textit{How do reference points contribute to understanding human welfare and decision-making?}
\end{tcolorbox}

%%%%%%%%%%%%%%%%%%%%%%%%%%%%%%%%%%%%%%%%%%%%%%%%%%%%%%%%%%%%%%%%%%%%%%%%%%%%%%%
% CHAPTER OVERVIEW
%%%%%%%%%%%%%%%%%%%%%%%%%%%%%%%%%%%%%%%%%%%%%%%%%%%%%%%%%%%%%%%%%%%%%%%%%%%%%%%
\begin{tcolorbox}[
    colback=green!5!white,
    colframe=green!75!black,
    title={\textbf{Chapter Overview}},
    fonttitle=\bfseries
]
This chapter develops the concept of the reference structure $C^*$ through five sections:

\begin{enumerate}[nosep]
    \item \textbf{Section \ref{sec:cstar-intro}}: Why $C^*$ matters---motivation and the Stigler-Becker parallel
    \item \textbf{Section \ref{sec:cstar-derivations}}: Three independent derivations (axiomatic, evolutionary, empirical)
    \item \textbf{Section \ref{sec:cstar-calibration}}: Calibration requirements and limitations
    \item \textbf{Section \ref{sec:cstar-integration}}: Integration with the 10C CORE structure
    \item \textbf{Section \ref{sec:cstar-summary}}: Summary and key takeaways
\end{enumerate}
\end{tcolorbox}

%==============================================================================
% DIE 9 FRAGEN NAVIGATION BOX
%==============================================================================
\BCMQuestionNav{6}

%==============================================================================
% 10C CORE CONNECTION BOX
%==============================================================================
\begin{tcolorbox}[
    colback=blue!5!white,
    colframe=blue!75!black,
    title={\textbf{10C CORE Connection: Foundation for HOW}},
    fonttitle=\bfseries
]
\textbf{Foundation Question:} What is the reference structure $C^*$?

\medskip
\begin{tabular}{ll}
\textbf{CORE Appendix:} & D (FORMAL-CSTAR) \\
\textbf{Output:} & Self-consistent structure $C^* = \Phi(C^*, \Psi)$ \\
\textbf{Key Insight:} & $C^*$ is derived, not assumed \\
\textbf{Novel Contribution:} & Three convergent derivations
\end{tabular}

\medskip
\textbf{The Core Equation:}
\begin{center}
$C^* = \Phi(C^*, \Psi)$ \quad (Fixed-point: beliefs reproduce themselves)
\end{center}
\end{tcolorbox}

%==============================================================================
% APPENDIX REFERENCES BOX
%==============================================================================
\begin{tcolorbox}[
    colback=yellow!10!white,
    colframe=orange!75!black,
    title={\textbf{Appendix References for This Chapter}},
    fonttitle=\bfseries
]
\begin{tabular}{lll}
\textbf{Appendix} & \textbf{Content} & \textbf{Section} \\
\midrule
D (FORMAL-CSTAR) & Complete C* theory: axioms, theorems, proofs & All \\
B (CORE-HOW) & Complementarity structure $\Gamma$ & \S\ref{sec:cstar-derivations} \\
BBB (CORE-WHERE) & Calibration of $C^*(\Psi)$ & \S\ref{sec:cstar-calibration}
\end{tabular}

\medskip
\textbf{Reading Guide:} This chapter provides intuition; Appendix UNMAPPED_PRF contains
the complete mathematical theory with all proofs.
\end{tcolorbox}


%%%%%%%%%%%%%%%%%%%%%%%%%%%%%%%%%%%%%%%%%%%%%%%%%%%%%%%%%%%%%%%%%%%%%%%%%%%%%%%
%                    PART I: INTRODUCTION
%%%%%%%%%%%%%%%%%%%%%%%%%%%%%%%%%%%%%%%%%%%%%%%%%%%%%%%%%%%%%%%%%%%%%%%%%%%%%%%

\section{Why $C^*$ Matters}
\label{sec:cstar-intro}

The coherence index $K$ measures deviation from a reference structure $C^*$:
\begin{equation}
\label{eq:coherence-index}
K = 1 - \frac{\|C - C^*\|}{\|C^*\|}
\end{equation}

For this to be scientifically meaningful, $C^*$ must be \emph{derived}, not assumed.

\subsection{The Trust Example}
\label{sec:ch6-trust-example}

\begin{intuition}[The Trust Example]
Merchants deciding whether to extend credit face three scenarios:
\begin{itemize}[nosep]
    \item \textbf{Beliefs match reality:} Trust is rewarded---beliefs reproduce
    \item \textbf{Too optimistic:} Get cheated, revise downward
    \item \textbf{Too pessimistic:} Miss opportunities, discover cooperation works
\end{itemize}
In all cases, beliefs gravitate toward $C^*$---the structure that reproduces itself.
\end{intuition}

\textbf{Numerical Illustration:} Consider Maria, a merchant with initial trust $C_0 = 0.9$
(highly optimistic). After 5 interactions where 30\% of partners defect, her beliefs
update: $C_1 = 0.82$, $C_2 = 0.76$, $C_3 = 0.73$, converging toward $C^* \approx 0.72$.

\subsection{The Stigler-Becker Parallel}
\label{sec:ch6-stigler-becker}

\citet{stiglerbecker1977} argued that economists should assume stable preferences.
We preserve this insight at a higher level: not stable preferences $U$, but stable
reference structure $C^*(\Psi)$. The key advance: $C^*$ is not exogenously assumed
but derived as the unique self-consistent configuration.


%%%%%%%%%%%%%%%%%%%%%%%%%%%%%%%%%%%%%%%%%%%%%%%%%%%%%%%%%%%%%%%%%%%%%%%%%%%%%%%
%                    PART II: THREE DERIVATIONS
%%%%%%%%%%%%%%%%%%%%%%%%%%%%%%%%%%%%%%%%%%%%%%%%%%%%%%%%%%%%%%%%%%%%%%%%%%%%%%%

\section{Three Derivations of $C^*$}
\label{sec:cstar-derivations}

Appendix UNMAPPED_PRF provides three independent derivations that all yield the same $C^*$.

\subsection{Axiomatic Derivation (Self-Consistency)}
\label{sec:ch6-axiomatic}

$C^*$ is the unique fixed point satisfying:
\begin{equation}
\label{eq:cstar-fixedpoint}
C^* = \Phi(C^*, \Psi)
\end{equation}

Under regularity, reciprocity, and contraction axioms (D-REG, D-REC, D-CON),
the Banach fixed-point theorem guarantees existence and uniqueness.

\textbf{Numerical Example:} Consider a trust game with $\alpha = 0.6$, $\beta = 0.3$,
and context sensitivity $\gamma = 0.15$. Starting from any initial $C_0$, iterating
$C_{n+1} = \Phi(C_n, \Psi)$ converges to $C^* \approx 0.72$ within 8 iterations.
\emph{See Appendix UNMAPPED_PRF, Part II.}

\subsection{Evolutionary Derivation (Stability)}
\label{sec:ch6-evolutionary}

$C^*$ is the evolutionarily stable structure (ESS). No mutant belief $C' \neq C^*$
can invade a population at $C^*$. The stability condition requires:
\begin{equation}
\label{eq:cstar-ess}
\pi(C^*, C^*) > \pi(C', C^*) \quad \forall C' \neq C^*
\end{equation}
where $\pi(C, C')$ denotes the fitness of strategy $C$ against $C'$.

Naive cooperators ($C = 1.0$) get exploited; paranoid defectors ($C = 0.0$)
miss cooperative gains. Both are outcompeted by accurate beliefs ($C = C^* \approx 0.72$).
\emph{See Appendix UNMAPPED_PRF, Part III.}

\subsection{Empirical Derivation (Long-Run Average)}
\label{sec:ch6-empirical}

$C^*$ is the ergodic limit:
\begin{equation}
\label{eq:cstar-ergodic}
C^*_{emp} = \lim_{T\to\infty} \frac{1}{T}\sum_{t=1}^{T} C_t
\end{equation}

Under adaptive dynamics with unique stationary distribution, the long-run
average converges to the theoretical $C^*$.
\emph{See Appendix UNMAPPED_PRF, Part IV.}

\begin{theorem}[Convergence---Appendix UNMAPPED_PRF, Theorem D.5]
\label{thm:cstar-convergence}
Under the foundational axioms and ergodicity:
\begin{equation}
\label{eq:cstar-convergence}
C^*_{axiomatic} = C^*_{evolutionary} = C^*_{empirical}
\end{equation}
\end{theorem}


%%%%%%%%%%%%%%%%%%%%%%%%%%%%%%%%%%%%%%%%%%%%%%%%%%%%%%%%%%%%%%%%%%%%%%%%%%%%%%%
%                    PART III: CALIBRATION LIMITS
%%%%%%%%%%%%%%%%%%%%%%%%%%%%%%%%%%%%%%%%%%%%%%%%%%%%%%%%%%%%%%%%%%%%%%%%%%%%%%%

\section{Calibration Requirements}
\label{sec:cstar-calibration}

$C^*$ cannot be learned from text alone. Each derivation requires empirical input.

\subsection{Parameter Requirements by Derivation}
\label{sec:ch6-param-requirements}

Each derivation pathway has specific calibration needs:

\begin{itemize}[nosep]
    \item \textbf{Axiomatic:} Parameters $\alpha, \beta, \gamma$ (magnitudes not in axioms)
    \item \textbf{Evolutionary:} Fitness landscape (selection pressures not documented)
    \item \textbf{Empirical:} Behavioral observations (effect sizes scattered)
\end{itemize}

\subsection{Calibration Sources}
\label{sec:ch6-calibration-sources}

The EBF framework provides two main calibration pathways:
\begin{enumerate}[nosep]
    \item \textbf{Literature-based:} Appendix UNMAPPED_BBB (CORE-WHERE) aggregates parameter values from 1,900+ papers
    \item \textbf{LLM Monte Carlo:} Appendix UNMAPPED_LLM (METHOD-LLMMC) estimates parameters via LLM-based simulation
\end{enumerate}

\textbf{Typical Values:} For trust contexts, $\alpha \in [0.5, 0.8]$, $\beta \in [0.2, 0.4]$,
$\gamma \in [0.1, 0.2]$, yielding $C^* \in [0.65, 0.80]$.


%%%%%%%%%%%%%%%%%%%%%%%%%%%%%%%%%%%%%%%%%%%%%%%%%%%%%%%%%%%%%%%%%%%%%%%%%%%%%%%
%                    PART IV: 10C INTEGRATION
%%%%%%%%%%%%%%%%%%%%%%%%%%%%%%%%%%%%%%%%%%%%%%%%%%%%%%%%%%%%%%%%%%%%%%%%%%%%%%%

\section{10C Integration}
\label{sec:cstar-integration}

\begin{table}[htbp]
\centering
\caption{Chapter 6 Integration with 10C CORE Structure}
\label{tab:ch6-9c-integration}
\small
\begin{tabular}{@{}llp{5cm}l@{}}
\toprule
\textbf{CORE} & \textbf{Question} & \textbf{Connection} & \textbf{Ref} \\
\midrule
\cellcolor{purple!20}WHO & Who has utility? & Agents in $C^*$ computation & AAA \\
\cellcolor{red!20}WHAT & What is utility? & Dimensions in $C^*$ matrix & C \\
\cellcolor{blue!20}HOW & How interact? & $C^*$ is the stable $\Gamma$-matrix & B \\
\cellcolor{green!20}WHEN & Context? & $C^* = C^*(\Psi)$ depends on context & V \\
\cellcolor{gray!20}WHERE & Parameters? & Calibration of $C^*$ & BBB \\
\cellcolor{orange!20}AWARE & Awareness? & Beliefs about $C^*$ matter & AU \\
\cellcolor{violet!20}READY & Action threshold? & $C^*$ determines equilibrium & AV \\
\cellcolor{teal!20}STAGE & Journey phase? & $C^*$ defines stable configs & AW \\
\bottomrule
\end{tabular}
\end{table}


%%%%%%%%%%%%%%%%%%%%%%%%%%%%%%%%%%%%%%%%%%%%%%%%%%%%%%%%%%%%%%%%%%%%%%%%%%%%%%%
%                    PART V: SUMMARY
%%%%%%%%%%%%%%%%%%%%%%%%%%%%%%%%%%%%%%%%%%%%%%%%%%%%%%%%%%%%%%%%%%%%%%%%%%%%%%%

\section{Summary}
\label{sec:cstar-summary}

\begin{tcolorbox}[
    colback=blue!5!white,
    colframe=blue!75!black,
    title={\textbf{Chapter 6 Summary}}
]
\textbf{Core Insight:} $C^*$ is the unique self-consistent reference structure.

\textbf{Three Derivations} (all in Appendix UNMAPPED_PRF):
\begin{enumerate}[nosep]
    \item Axiomatic: Fixed-point $C^* = \Phi(C^*, \Psi)$
    \item Evolutionary: ESS condition
    \item Empirical: Ergodic long-run average
\end{enumerate}

\textbf{Key Results:}
\begin{itemize}[nosep]
    \item $C^*$ exists, is unique, and is context-dependent
    \item Three derivations converge on the same $C^*$
    \item $C^*$ requires calibration---cannot be learned from text
\end{itemize}

\textbf{Transition:} Chapter 7 uses $C^*$ to define coherence dynamics $K(t)$.
\end{tcolorbox}

%------------------------------------------------------------------------------
% FORMAL DETAILS
%------------------------------------------------------------------------------
\subsection{Formal Details}
\label{sec:ch6-formal-details}

The complete formal treatment appears in Appendix D (FORMAL-CSTAR). Key results:

\begin{itemize}[nosep]
    \item \textbf{Existence:} Under axiom D-REG (regularity), $\Phi$ maps compact sets to compact sets
    \item \textbf{Uniqueness:} Under axiom D-CON (contraction), $\|\Phi(C_1,\Psi) - \Phi(C_2,\Psi)\| \leq k\|C_1 - C_2\|$ with $k < 1$
    \item \textbf{Context Dependence:} $C^*(\Psi_1) \neq C^*(\Psi_2)$ for different contexts (Theorem D.3)
    \item \textbf{Convergence Rate:} $\|C_n - C^*\| \leq k^n \|C_0 - C^*\|$ (exponential convergence)
\end{itemize}

%------------------------------------------------------------------------------
% WHAT COMES NEXT
%------------------------------------------------------------------------------
\subsection{What Comes Next}
\label{sec:ch6-next}

With $C^*$ established, subsequent chapters build:

\begin{enumerate}[nosep]
    \item \textbf{Chapter 7:} Uses $C^*$ to define the coherence index $K = 1 - \|C - C^*\|/\|C^*\|$
    \item \textbf{Chapter 8:} Establishes the mathematical foundations for dynamics
    \item \textbf{Chapter 9:} Shows how context $\Psi$ modulates $C^*(\Psi)$
\end{enumerate}

%------------------------------------------------------------------------------
% READING PATH BOX
%------------------------------------------------------------------------------
\begin{tcolorbox}[
    colback=gray!10!white,
    colframe=gray!75!black,
    title={\textbf{Reading Paths}}
]
\textbf{For Complete Theory:} $\rightarrow$ Appendix UNMAPPED_PRF (FORMAL-CSTAR)

\textbf{For Practitioners:} $\rightarrow$ Chapter 7 (coherence) $\rightarrow$ Chapter 9 (context)

\textbf{For Calibration:} $\rightarrow$ Appendix UNMAPPED_EST $\rightarrow$ Appendix LLM1
\end{tcolorbox}

%%%%%%%%%%%%%%%%%%%%%%%%%%%%%%%%%%%%%%%%%%%%%%%%%%%%%%%%%%%%%%%%%%%%%%%%%%%%%%%
